\chapter{Evaluation and Testing}

This chapter describes the metrics and procedures used to evaluate the effectiveness, efficiency, and usability of the proposed system. The evaluation is conducted at three levels: retriever, generator, and overall system performance.

\section{Evaluation Metrics}

\subsection{Retriever Metrics}

\textbf{Recall@K} measures how many relevant documents are successfully retrieved within the top $K$ results for a given query:
\[
\text{Recall@K} = 
\frac{\lvert \text{Relevant Documents} \cap \text{Top-}K \text{ Retrieved Documents} \rvert}
{\lvert \text{Relevant Documents} \rvert}
\]
This metric reflects the ability of the retriever to cover all relevant information.

\textbf{Precision@K} evaluates the proportion of retrieved documents in the top $K$ that are actually relevant:
\[
\text{Precision@K} =
\frac{\lvert \text{Relevant Documents} \cap \text{Top-}K \text{ Retrieved Documents} \rvert}
{\lvert \text{Top-}K \text{ Retrieved Documents} \rvert}
\]
It indicates how accurate the retrieved results are.

\textbf{Retrieval Latency} represents the time required for the retriever to return results:
\[
\text{Latency (ms)} = \text{Result Returned Time} - \text{Query Submission Time}
\]

\textbf{Coverage} measures how often the system retrieves at least one relevant document across all queries:
\[
\text{Coverage} =
\frac{\text{Number of Queries with Relevant Results}}
{\text{Total Number of Queries}} \times 100
\]

\subsection{Generator Metrics}

\textbf{ROUGE Score} quantifies textual similarity between generated outputs and reference answers. For ROUGE-N:
\[
\text{ROUGE-N} =
\frac{
\sum_{\text{ref}} \sum_{\text{ngram} \in \text{generated}} \text{Count}(\text{ngram})
}{
\sum_{\text{ref}} \sum_{\text{ngram} \in \text{reference}} \text{Count}(\text{ngram})
}
\]

\textbf{Relevance} is assessed through human judgment, rating how well generated responses align with user intent. The final score is computed as:
\[
\text{Relevance} =
\frac{1}{N} \sum_{i=1}^{N} \text{Relevance Score}_i
\]

\textbf{Factual Accuracy} measures the proportion of generated responses that are factually correct based on retrieved evidence:
\[
\text{Factual Accuracy} =
\frac{\text{Number of Correct Outputs}}
{\text{Total Number of Outputs}} \times 100
\]

\textbf{Generation Latency} captures the time taken by the generator to produce an output after retrieval completes:
\[
\text{Latency (ms)} = \text{Output Generation Time} - \text{Retrieval Completion Time}
\]

\subsection{System-Level Metrics}

\textbf{Usability} is evaluated via user feedback, typically collected using a Likert scale. The overall usability score is averaged across tasks:
\[
\text{Usability} =
\frac{1}{N} \sum_{i=1}^{N} \text{User Rating}_i
\]

\textbf{Task Success Rate} measures the percentage of tasks that users complete successfully:
\[
\text{Task Success Rate} =
\frac{\text{Number of Completed Tasks}}
{\text{Total Number of Tasks}} \times 100
\]

\textbf{End-to-End Latency} reflects the total response time of the system from user input to final output delivery:
\[
\text{Integration Latency (ms)} =
\text{Output Delivery Time} - \text{Query Submission Time}
\]

\section{Evaluation Methodology}

\subsection{Retriever Evaluation}

The retriever is evaluated using a benchmark dataset with predefined relevance judgments. Performance is analyzed across different query categories:
\begin{itemize}
    \item \textbf{Simple Queries}: Direct queries requiring straightforward document retrieval.
    \item \textbf{Complex Queries}: Queries involving multiple constraints or ambiguous conditions.
    \item \textbf{Context-Aware Queries}: Queries that depend on contextual knowledge of the repository.
\end{itemize}

\subsection{Generator Evaluation}

The generator receives outputs from the retriever and is evaluated using:
\begin{itemize}
    \item \textbf{Automatic Metrics}: ROUGE scores to measure overlap with reference answers.
    \item \textbf{Human Evaluation}: Expert assessments of relevance, correctness, and response quality.
\end{itemize}
Special attention is given to challenging cases, such as incomplete or ambiguous retrieved information, to test robustness.

\subsection{System Evaluation}

System-level evaluation is conducted through controlled user studies where participants interact with the VSCode extension to perform predefined programming tasks. Feedback is collected through questionnaires and interviews to assess usability, responsiveness, and overall user experience.

\newpage